\section{Conclusions}
\label{sec:conclusions}

This study used 2,903 agent-based simulation runs to examine how source imitation, memory, timing, and WOM shape fake-news dissemination. The validated baseline reproduced configured false-publication rates and showed that truth-sensitive WOM reduces actual fake reposting relative to discovery-only WOM. The scenario analysis then demonstrated that source popularity is not a sufficient misinformation metric. Full attribute copying maximised the target source's selection gain but reduced false reposts because credibility also changed the target's publication process. Non-credibility copying preserved the target's high false-news propensity and provided the more defensible camouflage treatment.

The most concerning confirmed combination acted in period 1, retained endorsements for 25 periods, copied every non-credibility attribute, penalised false recommendations, and ignored true recommendations. It increased final-window fake reposts by 2.89 percentage points relative to its matched control. The result is relational rather than a claim that false-news penalties are harmful: the baseline shows the protective role of a balanced $-1/+1$ truth-sensitive policy. Early interventions and longer memory generally strengthened scenario effects, while very short memory largely erased them.

The experiments also show the value of explicit controls and reproducible randomness. Network/WOM and camouflage/control contrasts answer different questions; common seeds reproduced condition rankings and usually improved precision. These practices help prevent attractive but misleading conclusions from a stochastic scenario grid.

FAKENEWS-ABM should be understood as a conditional theory made executable. Its present evidence identifies plausible mechanisms, clarifies measurement, and supports future hypotheses. Empirical calibration, a separate objective truth-generation parameter, homophily, and preregistered replication are required before using it for platform-specific prediction. Within those boundaries, the model shows that the riskiest imitation is not necessarily copying everything. It is copying the attributes that attract attention while preserving the underlying propensity to publish false news.
